%!TeX root=csp-audit.tex
\section{Claims examined and withdrawn}\label{sec:refuted}

Two claims were raised, survived a round of review, and turned out to be
wrong. They are the most instructive entries in this document, because both
failed in the same way and it is the way a defect list is prone to fail:
reading the formal-looking part of a paper and not the sentence before it.

\subsection{The claimed citation cycle in \S5}\label{sec:refuted-cycle}

The claim was a cycle
$\text{Lemma }9\Rightarrow7\Rightarrow86\Rightarrow85\Rightarrow
\text{Cor }22\Rightarrow\text{Thm }21\Rightarrow\text{Lemma }8\Rightarrow
\text{Lemma }9$, closed at \texttt{SS:1214}, and that it ``has to be broken
before anything in \S5.4 or \S5.7 can be formalized''.

\textbf{Refuted mechanically.} Every \verb|\begin{proof}| in \texttt{main.tex},
\texttt{StrongSubalgebras.tex} and \texttt{necessary\brk Claims.tex} was parsed,
attributed to its statement (including the \verb|\newtheorem*| duplicates the
paper uses to restate results proved later), and its \verb|\ref|s extracted.
\textbf{58 statements with proofs, 152 edges, no cycles.}

The named edge is real: \texttt{LEMBlockOfGoodBridgeDoesNotHaveBAC} does invoke
the relational stable-intersection corollary at \texttt{SS:1214}, in its
$T=\TC$ case. But it closes nothing. That corollary's transitive proof-closure
is 25 statements and does not contain
\texttt{LEMBlockOfGoodBridgeDoesNotHaveBAC}; its direct dependencies are
Propagation and the stable-intersection theorem, and the latter's closure is 15
statements, also disjoint from it.

Consequence: \S5 can be written in dependency order, and there is no hidden
simultaneous induction to untangle. The claim would not have survived ten
minutes of parsing.

\subsection{The recursion-depth bound}\label{sec:refuted-depth}

The claim was that the depth bound $|A|+|\Gamma|$ of \cite[Lem.~5.2]{ZhukJACM}
was ``not established for the \textsc{CheckTuple} $\to$ \textsc{SolveNonlinked}
path'', and that ``without the repair the algorithm is not polynomial''.
\textbf{Withdrawn.}

Lemma~5.2 bounds the depth by arguing that every recursion site except
\textsc{CheckWeakerInstance} ``reduce[s] all domains of size greater than 1''.
At the \textsc{SolveNonlinked} site this needs the instance to be neither
linked nor fragmented --- otherwise the linked-component congruence can be
trivial on some domain, that domain does not shrink, and the bound degrades
from the constant $|A|$ to $O(|A|\cdot n)$, which is fatal, since the total
work is $\mathrm{poly}(N)^{\text{depth}}$. The source does state the hypothesis
and does discharge it. It is simply not in one place:

\begin{itemize}\tightlist
\item the precondition ``an instance that is not linked and not fragmented'' is
  in the prose introducing the function (\texttt{1704.01914:892}), and
  \emph{not} in the \texttt{Input:} line of its pseudocode;
\item it is discharged at the call site (\texttt{:635}) by ``$\Theta_{0}$ is
  not fragmented because it is crucial'' --- itself a one-line argument the
  source does not give: a fragmented unsatisfiable instance has an
  unsatisfiable part, and no constraint of the \emph{other} part is then
  crucial;
\item the step from ``not fragmented'' to ``\emph{every} domain splits'' ---
  propagate the partition along a path, which exists by non-fragmentedness and
  is value-preserving by $1$-consistency --- appears once, at \texttt{:1070},
  inside the proof of a \emph{different} lemma, and is never cited here.
\end{itemize}

So the mathematics is present and the theorem is not in doubt. What is absent
is any single place where the six recursion sites of Lemma~5.2 are discharged.
The same pattern holds for the $\Gamma$-order argument, whose case (3) needs
the projections to be unchanged under weakening, which holds because the
instance is $1$-consistent when \textsc{CheckWeakerInstance} runs --- true, and
unstated.

\emph{Chasing it down did leave one genuine hole}, and closing it needed an
argument that is not in the source. \textsc{CheckTuple} calls
\textsc{SolveNonlinked} not on $\Theta_{0}$ but on $\Theta_{0}$ with domains
restricted, and restricting can only \emph{destroy} the ``not linked''
precondition.

\begin{lemma}[The \textsc{SolveNonlinked} call site]\label{lem:solvenonlinked-callsite}
Let $\Theta_{0}$ be cycle-consistent, non-fragmented and not linked, and let
$\Theta_{0}'$ be $\Theta_{0}$ with each domain restricted to a subset
$D_{x}'\subseteq D_{x}$. Then every domain of $\Theta_{0}$ of size greater than
one is strictly shrunk by the time \textsc{SolveNonlinked}$(\Theta_{0}')$
recurses into \textsc{Solve}.
\end{lemma}

\begin{proof}
Write $\tau_{x}=\mathrm{LinkedCon}(\Theta_{0},x)$, a congruence by
\cite[Lem.~6.2]{ZhukJACM}. First, $\tau_{x}$ is proper at \emph{every}
variable. Since $\Theta_{0}$ is not linked there are at least two connected
components of the graph on $\bigsqcup_{z}\{z\}\times D_{z}$; and every
component meets every variable, because from any $(x,a)$ non-fragmentedness
supplies a path to $y$ and $1$-consistency extends $a$ along it. So each
$D_{y}$ is partitioned into at least two nonempty $\tau_{y}$-blocks, each
therefore proper.

Now split on whether $\Theta_{0}'$ is linked.

\emph{Not linked.} Then \textsc{SolveNonlinked} recurses on a component, whose
trace on each $D_{y}$ is a $\tau_{y}$-block intersected with $D_{y}'$ --- a
proper subset of $D_{y}$ by the previous paragraph.

\emph{Linked.} Then for every $z$ and all $a,b\in D_{z}'$ some path of
$\Theta_{0}'$ connects them; every such path is a path of $\Theta_{0}$, so
$D_{z}'\times D_{z}'\subseteq\tau_{z}$, i.e.\ $D_{z}'$ lies inside a single
$\tau_{z}$-block. That block is proper, so $D_{z}'\subsetneq D_{z}$ --- every
domain has \emph{already} been shrunk, at the \textsc{CheckTuple} step, before
\textsc{SolveNonlinked} was called at all.
\end{proof}

\begin{entry}[Why this is the right argument]\label{ent:callsite-moral}
Two things are worth extracting.

First, the question this replaces was the wrong one. The natural worry is
whether the restriction preserves the stated precondition ``not linked'', and
one then starts asking how the minimal linear congruence $\mathrm{ConLin}$
relates to $\mathrm{LinkedCon}$ --- a relation the source never states, and
which is not needed. The dichotomy above uses no relation between them:
\emph{either} the precondition survives and the component split does the
shrinking, \emph{or} it fails and the restriction has already done it.

Second, it follows that \textsc{SolveNonlinked}'s precondition is not required
for the depth bound at all. The function is correct on a linked instance too
--- there is then one component and it recurses on the whole of it --- and the
lemma covers that case. The precondition should be read as documentation of the
intended use, not as an obligation the caller must discharge.
\end{entry}

\subsection{Declared moot: the arity in the special-WNU lemma}\label{sec:moot}

\cite[Lem.~4.7]{MarotiMcKenzie} is cited for the existence of a special
idempotent WNU of arity $n^{n!}$ in $\Clo(w)$. The concern raised was that the
period of $y\mapsto w(x,\dots,x,y)$ need not divide $n!$. The suggested
counterexample --- $n=3$ with a $5$-cycle --- does not work, since idempotence
makes $x$ a fixed point of that map and a full $5$-cycle on five elements is
then impossible; but cycles of length up to $|A|-1$ remain possible, and
$|A|-1$ need not divide $n!$ when $|A|>n$.

\emph{Moot.} The dichotomy proof consumes only ``there exist $N$ and a special
idempotent WNU of arity $N$ in $\Clo(w)$''. The precise arity matters only for
\S4 of \cite{ZhukSimplified}, which is an independent second result and out of
scope. The proof paper states the weaker form.

\subsection{A cycle in our own arrangement}\label{sec:our-cycle}

The parser that refuted the claimed cycle in the source
(\S\ref{sec:refuted-cycle}) was pointed at the proof paper, and found one
there:
\begin{center}\footnotesize
\pkey{cor:intersection-good} $\to$ \pkey{lem:propagation} $\to$
\pkey{lem:factor-by-delta} $\to$\\[2pt]
\pkey{lem:cut-or-nothing} $\to$ \pkey{lem:leftlinked-stay-full} $\to$\\[2pt]
\pkey{lem:center-pushed-in} $\to$ \pkey{cor:intersection-good}
\end{center}
Every edge was a real citation. The cycle was \emph{ours}: the source proves
the corollary from \texttt{LEMReverseHomomorphism} (\texttt{SS:953}), a
separate and much easier statement, where our draft had folded the backward
clauses into Propagation and cited the whole thing. Citing the bundle drags in
the forward clauses, which rest on \S5.6, which rests back on the corollary.

The fix was to state the backward clauses as their own lemma, which the
external review had independently asked for on different grounds --- it noted
that ``the reverse-homomorphism lemma used for (bt) is neither stated nor
cited''. The proof paper now has \pkey{lem:reverse-hom}, and its graph is
acyclic: $97$ statements, $52$ with proofs, $125$ edges.

Two things worth drawing out. First, the source was right and we were wrong,
which is the opposite of the direction this document usually runs. Second, this
is the third time the same tool has settled a question in minutes that prose
inspection had not settled at all.

\medskip

Fixing the first cycle exposed a second, and this one the review had predicted.
It wrote: ``There is also a possible dependency-cycle danger:
\texttt{block-good-bridge} invokes stable intersection; stable intersection
invokes \texttt{no-cross-bridge}; and the most natural proof of no-cross
appears to route through the linear bridge characterization and the block-good
machinery. An explicit acyclic dependency graph is needed.'' That is exactly
right, and once the two bridge-classification lemmas were written out the
parser closed the loop:
\begin{center}\footnotesize
\pkey{cor:stable-intersection} $\to$ \pkey{thm:stable-intersection} $\to$
\pkey{lem:bridge-to-pc} $\to$ \pkey{lem:no-cross-bridge} $\to$\\[2pt]
\pkey{lem:linear-equiv} $\to$ \pkey{lem:nontrivial-bridge} $\to$
\pkey{lem:block-good-bridge} $\to$ \pkey{cor:stable-intersection}
\end{center}
It is not vicious. \texttt{LEMBlockOfGoodBridgeDoesNotHaveBAC} appeals to
stable intersection at $n=3$ with all types $\TC$, and that case rests on
\cite[Lem.~6.11]{ZhukStrong} and on nothing in the bridge theory. The proof
paper therefore states that case as its own lemma, proved before the bridges,
and both consumers cite it.

The source avoids the loop by not citing: its
\texttt{LEMPCBridgesAreTrivial} contradicts PC-ness without naming
\texttt{LEMLinearEquivalentConditions}, so the edge that closes our cycle is
absent from its graph. The dependency is nevertheless there in the mathematics
--- ``$\sigma$ is PC, so no bridge from $\sigma$ to $\sigma$ has
$\bcol\delta\supsetneq\sigma$'' \emph{is} that lemma --- which is worth
recording, because it means the earlier acyclicity verdict on the source
(\S\ref{sec:refuted-cycle}) is a statement about its citations and not about
its logical dependencies. The claim refuted there was a specific named cycle,
and that refutation stands; the stronger reading, that no cycle of any kind is
latent in \S5, was never established and is now known to be delicate.

\section{The machine checks}\label{sec:machine-checks}

Everything reducible to a finite check has had one. Each entry gives the scope,
so that a reader can see what it does and does not cover. Two of the original
claims of this document died to exactly such checks.

\begin{center}\footnotesize
\begin{tabular}{@{}p{0.20\textwidth}p{0.30\textwidth}p{0.38\textwidth}@{}}
\hline
\textbf{Check} & \textbf{Scope} & \textbf{Result} \\
\hline
Proof-dependency graph, the source & 58 statements, 152 edges over the three
  \TeX{} files & Acyclic; \S\ref{sec:refuted-cycle} \\
Proof-dependency graph, the proof paper & 97 statements, 52 with proofs, 125
  edges & Acyclic --- after one real cycle was found and fixed,
  \S\ref{sec:our-cycle}. Run by the build gate \\
Statement audit & all 81 labelled statements & 9 flagged, 8 cleared by hand;
  $\mathrm D2$ is the sole survivor \\
$\mathrm C1$: special idempotent affine WNU on $\mathbb Z_{p}$ &
  $p\in\{2,3,5,7\}$, $3\le n<12$ & exists exactly when $p\mid n-1$, and is then
  the sum \\
$\mathrm D6$: the third hypothesis can fail &
  $\mathbb Z_{4}\times\mathbb Z_{2}\times\mathbb Z_{4}$ & witness found;
  \S\ref{sec:D6} \\
$\mathrm D{12}$: counterexample & all six bijections $\mathbb Z_{3}\to
  \mathbb Z_{3}$, exhaustive over $\mathbb Z_{3}^{4}$ & none carries $\delta$
  to the difference relation \\
$\mathrm D{14}$: the subrelation form &
  nine small Taylor algebras (min-semilattices, the dual discriminator,
  $\mathbb Z_{2}$ minority, $\mathbb Z_{3}$, $\mathbb Z_{2}^{2}$, a semilattice
  square), $527+99\,477$ pairs $(W,W')$ & no counterexample; and $0_{C}$
  absorbs $C^{2}$ in no algebra tested \\
\hline
\end{tabular}
\end{center}

\begin{entry}[What the searches are and are not]\label{ent:search-limits}
A search over nine small algebras is evidence that a statement is not
\emph{obviously} false, and nothing more; the two positive entries above
($\mathrm C1$ and $\mathrm D14$) are backed by written proofs in the proof
paper, and the searches were run first, to decide whether the proofs were worth
attempting. The two \emph{negative} entries are different in kind: a
counterexample found by exhaustion over a finite structure is a proof, and
$\mathrm D6$ and $\mathrm D12$ rest on nothing else. Those two are stated in
the proof paper in a form a reader can check by hand in a few lines, which is
the right standard for a paper; the exhaustive verification is a check on the
transcription, not the argument.

Where the proof paper still says ``verified by exhaustion'' for a claim with no
short symbolic proof, that is a gap, and it is listed in
\S\ref{sec:risk}.
\end{entry}

\section{Disposition of the external review}\label{sec:review}

An external reviewer read a single-file build of an earlier draft. Its
findings, and what was done with each.

\subsection{Accepted}\label{sec:review-accepted}

\begin{itemize}\tightlist
\item \textbf{$\mathrm D17$ is broken at its first line.} Correct, and it
  reverses a claim made here. \S\ref{sec:D17}; the item is reopened.
\item The symmetrisation lemma cited bridge composition, which carries an
  irreducibility hypothesis it does not have. Fixed by proving the collapse
  identity directly.
\item The nice-bridge lemma asserted that $\delta\ast\delta$ is an equivalence
  relation, which needs a stabilised power or the Maltsev route.
\item The block-good-bridge proof treated $(\omega\cup\omega^{-1})\cap D^{2}$ as
  a subalgebra, which a union is not; use $(\omega\circ\omega^{-1})\cap D^{2}$
  (\S\ref{sec:exp-54}).
\item ``Which we may take reflexive'' is used twice with no lemma behind it.
  Accepted as a gap in our exposition, but the gap is not where either of us
  thought: there is nothing to normalise. Every consumer hypothesises
  $\bcol\delta\supsetneq\sigma$, which contains $\sigma\subseteq\bcol\delta$,
  and $(a,a)\in\sigma\subseteq\bcol\delta$ \emph{is} $(a,a,a,a)\in\delta$. So
  the bridge is reflexive already, and the proof paper says so in one line
  (\pkey{lem:bridge-reflexive}). A genuine normalisation of the same
  \emph{shape} does survive elsewhere --- ``we may assume $\bcol\delta$
  rectangular, as otherwise compose it with itself many times'', in
  \texttt{LEMNoBridgeBetweenDifferentTypes} --- and that one is open.
\item The abelian-bridge characterisation needs $|A|>1$, and ``composing enough
  times'' needs an explicit stabilised power.
\item The reason given for reflexivity in \pkey{lem:perfect-from-linked} is
  false: $(a,a)$ lying in the first projection gives some $\delta(a,a,c,d)$,
  not $\delta(a,a,a,a)$. The repair is the reviewer's: clause (d) gives
  $c\,\sigma\,d$ and stability replaces $d$ by $c$.
\item The document had the wrong shape, mixing a proof paper with a project
  status report. That is what this three-way split answers.
\item A duplicate label, several overfull boxes, a \verb|\ref| printed inside a
  \verb|\cite|, and an unidentifiable bibliography entry. All fixed; the build
  gate now fails on each class.
\end{itemize}

Several further findings concern statements the proof paper had not yet written
in full --- the induction measure across coverings (Entry~\ref{ent:measure}),
weakening's termination order, the common binary absorption witness
(\S\ref{sec:D1}), the mixed-prime linear algebra, and the definition of
$\textsc{Solve}$. They are scheduled work rather than audit entries, and they
are listed in the proof paper's status table.

\subsection{Accepted, with a correction}\label{sec:review-corrected}

The reviewer wrote that $\LeftLinked$ and $\RightLinked$ ``are used throughout
the strong-subalgebra proofs \dots but no definition is given in the
manuscript, and Brady's notes do not use those exact names''. Both halves are
true of \emph{our} draft and of \cite{BradyNotes}. But
\cite{ZhukSimplified} \textbf{does} define them, at \texttt{SS:52--66}, and in
the general $n$-ary form:

\begin{quote}\small
``For a relation $R\le\mathbf A_{1}\times\dots\mathbf A_{n}$ by
$\LeftLinked(R)$ we denote the minimal equivalence relation on $\proj_{1}(R)$
such that $(a_{1},a_{2},\dots,a_{n}),(b_{1},a_{2},\dots,a_{n})\in R$ implies
$(a_{1},b_{1})\in\LeftLinked(R)$. Similarly, $\RightLinked(R)$ is the minimal
equivalence relation on $\proj_{n}(R)$ such that
$(a_{1},\dots,a_{n-1},a_{n}),(a_{1},\dots,a_{n-1},b_{n})\in R$ implies
$(a_{n},b_{n})\in\RightLinked(R)$.''
\end{quote}

So this is not a defect of the source but an omission of ours: we reproduced
the uses without the definition. It is now
\pkey{def:linked} and \pkey{lem:linked-basics}, with the equivalence to
connectedness of the bipartite graph proved for subdirect binary relations.

The $n$-ary form settles a question that matters elsewhere. $\LeftLinked$
splits $R$ at its \emph{first} coordinate and $\RightLinked$ at its
\emph{last}, with all intermediate coordinates going to the other side --- so
$\RightLinked(\proj_{1,2,3}(\delta))$ concerns the split $(1,2)\mid(3)$ and is
a relation on the third coordinate. That is exactly the reading $\mathrm D14$
(\S\ref{sec:D14}) turns on, and it is now confirmed against the definition
rather than inferred from the uses.

\subsection{Rejected, with reasons}\label{sec:review-rejected}

\begin{itemize}\tightlist
\item \textbf{The witness in \texttt{LEMBridgeBetweenCongruences} is correct as
  printed.} The reviewer proposed reading $z_{i}=y_{i}=x_{i}$ for
  $z_{i}=y_{i}=x_{1}$. The printed choice works: it gives
  $(x_{1},x_{2},x_{1},x_{1})\in\delta$ using symmetry of $\sigma_{1}$ and
  reflexivity of $\omega$, which is what is needed for
  $\sigma_{1}\subseteq\proj_{1,2}(\delta)$. The proposed choice needs
  $(x_{1},x_{2})\in\omega$, which is not available.
\item \textbf{$\subt{\TBA,\TC}$ is a conjunction, not a disjunction.}
  \texttt{main:1567} reads ``we write $C<_{\TBA,\TC}B$ meaning that
  $C<_{\TBA}B$ and $C<_{\TC}B$''. The reviewer is right that the notation was
  never defined in our document, and it now is; but the reading is ``both''.
\end{itemize}

\section{Risk register}\label{sec:risk}

What is least well tested, stated so that a later reader knows where to push.

\begin{description}\tightlist
\item[Open mathematics.] $\mathrm D17$ (\S\ref{sec:D17}); the
  rectangularity normalisation in
  \texttt{LEMNoBridge\brk Between\brk DifferentTypes}; the induction
  measure across expanded coverings (Entry~\ref{ent:measure}). These gate the
  statements that depend on them and are marked \emph{open} in the proof
  paper's status table. The reflexivisation, listed here in an earlier draft,
  turned out to need no argument at all.

\item[Least-tested repairs.] $\mathrm D7$(ii)'s reweakening and $\mathrm D6$'s
  cruciality corollary each touch the CSP-instance machinery rather than the
  algebra, and each has been checked against only the call site that motivated
  it.

\item[Least-tested positive verdicts.] The self-intersection lemma was read only
  where the central-relation call sites forced it open. \S5.3's three small
  lemmas were opened last, and one of them is what splits the main case of the
  largest proof in \S5.

\item[Unread.] \texttt{main.tex}'s CSP half --- instances, coverings,
  cruciality, the main induction, the algorithm --- has been read only where
  $\mathrm D6$, $\mathrm D7$ and $\mathrm D8$ touch it. That is roughly two
  thousand lines to which the standard of \S\ref{sec:standard} has not been
  applied, and the register above should not be read as covering them. \S4 of
  the source (\texttt{XYSymmetric.tex}, 1\,640 lines) is out of scope by
  decision.

\item[Claims resting on exhaustion alone.] None of the register's positive
  claims, but the proof paper still carries ``verified by exhaustion'' where a
  short symbolic argument would serve --- for instance the $\mathbb Z_{3}$
  counterexample of $\mathrm D12$, where setting
  $(x_{1},x_{2},x_{3},x_{4})=(0,b,b,0)$ forces $2(f(0)-f(b))=0$, so $f$ is
  constant for odd $p$ and no bijection can work. Replace them.
\end{description}

\begin{entry}[The failure mode this document is prone to]\label{ent:meta}
Three times now a claim here has been published, reviewed, and turned out to be
wrong: the citation cycle, the recursion-depth bound, and the repair of
$\mathrm D17$. In each case the error survived internal review and died to
someone reading the argument from outside. A list like this one is exactly as
reliable as the reading behind each entry, and the failure mode is not
carelessness --- it is confidence in a step that was never read at the standard
of \S\ref{sec:standard}. The entries most at risk are the ones in the
\emph{unread} row above.
\end{entry}
